\section{Joining finite histories and determining the critical trajectory}
\label{sec:trajectory}

The bulk estimates provide a small trajectory on each finite scale
window.  The histories have different entry scales, so they need not
agree as activity families on their overlaps.  Their physical readouts
nevertheless determine the two marginal coordinates accurately enough
to compare them.  We first use those comparisons to form one infinite
sequence, and then determine its leading coefficient.  The coefficient
of the gradient activity obtained at the end of this section is the
one needed in the magnetic terminal calculation.

\subsection{An abstract lemma for sparse changes of history}

An ordinary evolution step increases the scale index by one.  A change
of history takes place at a fixed scale index and does not count as an
additional evolution step.  The following notation keeps these two
operations distinct.  We write $(u_j,Z_j)$ for the state used to start
the ordinary step at index $j$, and
$(\widetilde u_{j+1},\widetilde Z_{j+1})$ for the output of that step.
At an index where no change of history occurs, the tilded state is
$(u_{j+1},Z_{j+1})$.  At a change of history, it is the state immediately
before the change.

\begin{lemma}[A trajectory with geometrically separated jumps]
\label{lem:trajectory-sparse-splice}
Let $H,b>0$, let $j_0\ge1$, and let $\mathcal Q$ be a set of integers
larger than $j_0$.  Suppose that the increasing enumeration $(q_r)$ of
$\mathcal Q$ satisfies
\[
 q_{r+1}\ge \rho q_r
 \qquad\text{for some fixed }\rho>1.
\]
The set $\mathcal Q$ may also be finite or empty.  For every integer
$j\ge j_0$, let $(u_j,Z_j)$ and
$(\widetilde u_{j+1},\widetilde Z_{j+1})$ be real states with positive
second coordinates.  Assume that both coordinates of every state tend
to zero as its scale index tends to infinity.  Suppose that the
ordinary evolution satisfies
\begin{align}
 \widetilde u_{j+1}-u_j
   &=-H_j Z_j^2
       +O\bigl((|u_j|+Z_j^2)Z_j^2\bigr),
       \label{eq:trajectory-ordinary-u}\\
 \frac{\widetilde Z_{j+1}}{Z_j}
   &=1-b_j u_j+O(u_j^2+Z_j^2),
       \label{eq:trajectory-ordinary-z}\\
 H_j&=H+O(j^{-2}),\qquad b_j=b+O(j^{-2}).
       \label{eq:trajectory-coefficients}
\end{align}
For $q>j_0$ with $q\notin\mathcal Q$, assume
$(u_q,Z_q)=(\widetilde u_q,\widetilde Z_q)$.  At every $q\in\mathcal Q$,
assume instead
\begin{align}
 |u_q-\widetilde u_q|&\le C\widetilde Z_q^2,
       \label{eq:trajectory-jump-u}\\
 |Z_q-\widetilde Z_q|
   &\le C\bigl(|\widetilde u_q|+\widetilde Z_q^2\bigr)
                  \widetilde Z_q.
       \label{eq:trajectory-jump-z}
\end{align}
All constants in these hypotheses are uniform in the scale index.
Then, with $x_*=(H/b)^{1/2}$,
\begin{align}
 \frac{u_j}{Z_j}&=x_*+O(j^{-1}),
       \label{eq:trajectory-ratio}\\
 \frac1{Z_j}&=b x_*j+O(\log j),
       \label{eq:trajectory-inverse}\\
 u_j&=\frac1{bj}+O\left(\frac{\log j}{j^2}\right),\qquad
 Z_j=\frac1{\sqrt{bH}\,j}
                 +O\left(\frac{\log j}{j^2}\right).
       \label{eq:trajectory-asymptotics}
\end{align}
The same estimates hold for the states immediately before jumps.
\end{lemma}

\begin{proof}
We first obtain positive, comparable coordinates without assuming a
rate of convergence.  We then derive the rate $j^{-1}$, and finally
identify its coefficient by solving the equation for $u/Z$ backwards.

\paragraph{Grouping a jump with preceding evolution.}
Fix an integer $K$, whose size will be chosen below.  Discard finitely
many initial indices so that the intervals $[q-K,q]$, $q\in\mathcal Q$,
are disjoint.  This is possible because the gaps between consecutive
elements of $\mathcal Q$ tend to infinity.  Replace the $K$ ordinary
steps from $q-K$ to the state immediately before $q$, together with
the jump at $q$, by one grouped transition.  Retain every other
ordinary step.  Thus consecutive retained indices have difference
either $1$ or $K$.  A grouped transition of length $K$ still contains
exactly $K$ ordinary evolution steps.

Write $(u,Z)$ for the state at the beginning of such a transition.
For fixed $K$, induction in the ordinary recurrence gives, at each
of its intermediate ordinary states,
\[
 Z_{\mathrm{int}}/Z=1+o(1),\qquad
 u_{\mathrm{int}}=u+O_K(Z^2).
\]
The errors are uniform along the tail.  More precisely, immediately
before the jump the $K$ ordinary steps give
\begin{align}
 \widetilde u_q-u
   &=-KHZ^2+O_K\bigl((|u|+Z^2+j^{-2})Z^2\bigr),
       \label{eq:trajectory-k-u}\\
 \widetilde Z_q/Z
   &=1-Kbu+O_K(u^2+Z^2+j^{-2}|u|),
       \label{eq:trajectory-k-z}
\end{align}
where $j=q-K$.  To see these estimates, use the first recurrence to
bound the change in $u$ by $O_K(Z^2)$, and substitute this bound in the
second recurrence.  Multiplying its $K$ ratios creates only terms
bounded by $O_K(u^2+Z^2)$.  Substitution back into the first recurrence
then gives \eqref{eq:trajectory-k-u}.

The jump changes the right side of \eqref{eq:trajectory-k-u} by at
most $(C+o(1))Z^2$.  Its contribution to
\eqref{eq:trajectory-k-z} has absolute value at most
$(C+o(1))|u|+O_K(Z^2)$.  Here the leading constant $C$ is a constant
from the jump hypotheses, independent of $K$; the terms caused by
$\widetilde u_q-u$ are included in $O_K(Z^2)$.  Any error bounded by
$C_1|u|+C_KZ^2$ can be written as a coefficient of $u$, of absolute
value at most $C_1$, plus an $O_K(Z^2)$ term.  This remains true at
$u=0$, when the first term is zero.

Choose $K$ sufficiently large compared with the jump constants and
with $H^{-1},b^{-1}$.  After moving farther along the tail, every
retained transition therefore satisfies
\begin{equation}
 u^+-u=-\alpha Z^2,\qquad
 \frac{Z^+}{Z}=1-\beta u+r,
 \qquad |r|\le C_K Z^2,
 \label{eq:trajectory-grouped}
\end{equation}
where
\begin{equation}
 0<\alpha_-\le\alpha\le\alpha_+<\infty,
 \qquad
 0<\beta_-\le\beta\le\beta_+<\infty.
 \label{eq:trajectory-grouped-bounds}
\end{equation}
The bounds also hold for a retained single ordinary step.  The
coefficients $\alpha,\beta$ may vary from one transition to another;
no regularity of that dependence is used.

\paragraph{Positivity and the two comparison cones.}
At retained nodes, $u$ strictly decreases, because $Z>0$.  Its limit
is zero by hypothesis, so $u>0$ at every node of this tail.  We next
show that this positive coordinate is comparable to $Z$.

Choose $c_->0$ so small that $c_-^2\beta_+<\alpha_-/2$.
Whenever $u\le c_-Z$, \eqref{eq:trajectory-grouped} gives
\[
 (u^+-c_-Z^+)-(u-c_-Z)
   =-\alpha Z^2+c_-\beta uZ-c_-rZ
   \le -\tfrac14\alpha_-Z^2
\]
on a sufficiently late tail.  Consequently the region
$u-c_-Z\le0$ is forward invariant, and once it is entered the
quantity $u-c_-Z$ becomes strictly negative and then decreases.
This contradicts its convergence to zero.  Thus $u>c_-Z$.

Choose $c_+>0$ so large that
$c_+^2\beta_->2\alpha_+$.  Whenever $u\ge c_+Z$, the same computation
gives
\[
 (u^+-c_+Z^+)-(u-c_+Z)
   =-\alpha Z^2+c_+\beta uZ-c_+rZ
   \ge hZ^2
\]
for some $h>0$ on a sufficiently late tail.  This region is also
forward invariant.  Entry would make $u-c_+Z$ strictly positive and
increasing, again contradicting convergence to zero.  We have proved
\begin{equation}
 c_-Z<u<c_+Z
 \label{eq:trajectory-cones}
\end{equation}
at all sufficiently late retained nodes.

\paragraph{The first rate estimate.}
Put $x=u/Z$ at a retained node.  From
\eqref{eq:trajectory-grouped} and \eqref{eq:trajectory-cones},
\[
 \frac1{Z^+}-\frac1Z
   =\frac{\beta u-r}{Z(1-\beta u+r)}
   =\beta x+O(Z).
\]
This increment is bounded above and bounded away from zero.  Its
sum is therefore comparable to the number of retained transitions.
The physical scale index increases by either $1$ or $K$ at each
such transition, so that number is comparable to the scale index.
It follows that $Z\asymp j^{-1}$ and $u\asymp j^{-1}$ at retained
nodes.  The intermediate estimates preceding
\eqref{eq:trajectory-k-u}, together with the jump bounds, extend
these comparisons to every ordinary state and every state before a
jump.  In particular,
\begin{equation}
 Z_j\asymp j^{-1},\qquad u_j\asymp j^{-1},\qquad
 0<c\le u_j/Z_j\le C,
 \label{eq:trajectory-first-rate}
\end{equation}
and the corresponding tilded bounds hold as well.

We now know the order of both coordinates.  The remainder of the
proof identifies their leading coefficient.  For this purpose we
return to the original ordinary steps and keep each jump as a
separate event at its actual scale index.

\paragraph{Backward determination of the ratio.}
For an ordinary step, division of
\eqref{eq:trajectory-ordinary-u} by
\eqref{eq:trajectory-ordinary-z}, using
\eqref{eq:trajectory-first-rate}, gives
\begin{align}
 \frac{\widetilde u_{j+1}}{\widetilde Z_{j+1}}
          -\frac{u_j}{Z_j}
   &=(b x_j^2-H)Z_j+O(j^{-2})\notag\\
   &=b(x_j+x_*)(x_j-x_*)Z_j+O(j^{-2}),
 \qquad x_j=\frac{u_j}{Z_j}.
 \label{eq:trajectory-ratio-step}
\end{align}
At a jump, direct division and
\eqref{eq:trajectory-jump-u}--\eqref{eq:trajectory-jump-z} yield
\begin{equation}
 \left|\frac{u_q}{Z_q}
            -\frac{\widetilde u_q}{\widetilde Z_q}\right|
       \le Cq^{-1}.
 \label{eq:trajectory-ratio-jump}
\end{equation}

List all ordinary steps and all jumps in their chronological order,
and put $y=x-x_*$ at every state in this list.  An ordinary step
has the form
\[
 y^+=(1+a)y+e,\qquad
 a=b(x+x_*)Z,\qquad c/j\le a\le C/j,
 \qquad |e|\le Cj^{-2}.
\]
A jump has the same form with $a=0$ and $|e|\le C/q$.
For two positions $s<t$ in this event list, let $P_{s,t}$ be the
product of the factors $1+a$ between them.  The factors are at least
one, and $P_{s,t}\to\infty$ as $t\to\infty$ for fixed $s$, because
there is one ordinary step at every integer scale and
$\sum_jj^{-1}=\infty$.  Iterating the equation and solving for the
earlier value gives
\[
 y_s=\frac{y_t}{P_{s,t}}
       -\sum_{r=s}^{t-1}\frac{e_r}{P_{s,r+1}}.
\]
The values $y_t$ are bounded by \eqref{eq:trajectory-first-rate}.
Letting $t\to\infty$, and bounding the denominators below by one,
therefore shows, for a state at scale $j$,
\[
 |y_s|\le
 C\sum_{k\ge j}k^{-2}
   +C\sum_{\substack{q\in\mathcal Q\\q\ge j}}q^{-1}
 \le \frac Cj.
\]
The last bound uses geometric separation: if $q$ is the first jump
at or after $j$, the remaining reciprocals sum to at most
$q^{-1}/(1-\rho^{-1})$.  This proves
\eqref{eq:trajectory-ratio} at every state, including those before
jumps.

\paragraph{Counting ordinary time.}
For an ordinary step we now have
\begin{align*}
 \frac1{\widetilde Z_{j+1}}-\frac1{Z_j}
   &=b_j\frac{u_j}{Z_j}+O(j^{-1})
     =b x_*+O(j^{-1}).
\end{align*}
At a jump, \eqref{eq:trajectory-jump-z} and
\eqref{eq:trajectory-first-rate} give
\[
 \left|\frac1{Z_q}-\frac1{\widetilde Z_q}\right|\le C.
\]
There are $O(\log j)$ jumps up to scale $j$.  Telescoping from a
fixed initial index to $j$ counts exactly one ordinary increment
for each increase of the scale index by one.  The sum of the
ordinary errors is $O(\log j)$, and the sum of jump errors has the
same bound.  Hence \eqref{eq:trajectory-inverse} holds.  Inverting
that estimate and using \eqref{eq:trajectory-ratio} proves
\eqref{eq:trajectory-asymptotics}, since
$x_*/(b x_*)=1/b$.  A final jump changes $u$ and $Z$ by $O(j^{-2})$,
so these conclusions also hold for the states before jumps.
\end{proof}

\subsection{Application to the bulk histories}

For a dyadic integer $J$, attach a superscript $(J)$ to the
coordinates of the history whose entry index is $i(J)$ and whose
last index is $8J$.  The entry estimate gives $i(J)=O(\log J)$.
Thus, for every fixed $c>0$, the interval $[cJ,8J]$ lies after entry
when $J$ is large.  The estimates below are uniform on such intervals,
with constants allowed to depend on $c$.

\begin{lemma}[Comparison at a common scale]
\label{lem:trajectory-overlap}
Consider two bulk histories with comparable window parameters $J,J'$.
At an index $j$ belonging to fixed positive-fraction windows of both
histories, write $(u,Z)$ and $(u',Z')$ for their balanced marginal
coordinates.  For all sufficiently large indices,
\begin{equation}
 |u'-u|\le CZ^2,
 \qquad
 |Z'-Z|\le C(|u|+Z^2)Z.
 \label{eq:trajectory-overlap}
\end{equation}
The constants are uniform when the comparison factors for the two
windows are fixed.
\end{lemma}

\begin{proof}
The two readouts in Proposition~\ref{prop:bulk-readouts} concern the
same microscopic observables at scale $j$.  Consequently their values
can be written using either history as
\begin{align*}
 U_j&=\frac a{1-u}+O(Z^2)+O(e^{-c_0\min(J,J')}),\\
 Y_j&=M_jZ+O\bigl((|u|+Z^2)Z\bigr)
                   +O(e^{-c_0\min(J,J')}),
\end{align*}
and likewise with primes.  Here $a>0$ is fixed, and $M_j$ is common
to the two histories and bounded above and below by positive
constants.  An exponentially small discrepancy in $M_j$ can be
included in the displayed error.

By Lemmas~\ref{lem:bulk-window-vanishing} and
\ref{lem:bulk-odd}, the coordinates tend uniformly to zero,
$Z,Z'>0$, and $-\log Z,-\log Z'=o(\min(J,J'))$.  The exponential
errors are therefore negligible relative to every fixed positive
power of either $Z$ or $Z'$.  The second readout first implies
$Z'\asymp Z$: its error relative to its own leading term tends to
zero on each history.  Comparing the first readout now gives
\[
 \left|\frac a{1-u'}-\frac a{1-u}\right|\le CZ^2.
\]
The derivative of $a/(1-u)$ is bounded below positively near zero,
so $|u'-u|\le CZ^2$.  Substitute this estimate and $Z'\asymp Z$
back into the second readout.  After division by the positive lower
bound on $M_j$, it gives the second bound in
\eqref{eq:trajectory-overlap}.
\end{proof}

\begin{proposition}[The bulk trajectory throughout an overlap window]
\label{prop:bulk-window}
Fix the block factor $L$ and the other parameters chosen for the
bulk maps.  For every sufficiently large dyadic $J$, uniformly in
the $J$-history at integer indices $J/2\le j\le3J$, its unbalanced
gradient coordinate $t_j$, fundamental coordinate $z_j$, and plane
remainder $\mathcal R_j$ satisfy
\begin{equation}
 t_j=\frac1{2j\log L}
                 +O\left(\frac{\log j}{j^2}\right),
 \qquad |z_j|=O(j^{-1}),
 \qquad \|\mathcal R_j\|=O(j^{-2}).
 \label{eq:trajectory-bulk-window}
\end{equation}
The remainder norm is the unbalanced plane activity norm.
\end{proposition}

\begin{proof}
Select a dyadic $J_0$ large enough for the preceding estimates, and
use the $J$-history on $J\le j<2J$, for
$J=J_0,2J_0,4J_0,\ldots$.  At the index $2J$, first evolve the
$J$-history through its ordinary step from $2J-1$ to $2J$.  Regard
that output as the state before the change, and then replace it by
the state of the $2J$-history at the same index.  Both states lie in
the available positive-fraction windows.  Lemma~\ref{lem:trajectory-overlap}
gives \eqref{eq:trajectory-jump-u}--\eqref{eq:trajectory-jump-z} for
this replacement.

Lemma~\ref{lem:bulk-window-vanishing} makes all selected states and
all states before changes tend to zero.  Lemma~\ref{lem:bulk-odd}
provides $Z>0$ and the balanced remainder estimates
\[
 \|B_{e,j}\|\le CZ_j^2,
 \qquad
 \|B_{o,j}\|\le C(|u_j|+Z_j^2)Z_j.
\]
Inserting them into the bulk recurrences, with the parity structure
of Lemma~\ref{lem:bulk-parity}, gives
\eqref{eq:trajectory-ordinary-u}--\eqref{eq:trajectory-coefficients}
with
\[
 H>0,\qquad b=2\log L.
\]
The jumps occur at $2J_0,4J_0,8J_0,\ldots$, so they satisfy the
geometric-separation hypothesis with $\rho=2$.
Lemma~\ref{lem:trajectory-sparse-splice} consequently gives on the
selected path
\begin{equation}
 u_j=\frac1{2j\log L}
          +O\left(\frac{\log j}{j^2}\right),
 \qquad Z_j=O(j^{-1}),
 \qquad \|B_{e,j}\|=O(j^{-2}),
 \qquad \|B_{o,j}\|=O(j^{-2}).
 \label{eq:trajectory-balanced-window}
\end{equation}
The displayed odd bound is sufficient here; its stated forcing
estimate in fact gives $O(j^{-2})$ directly from $u_j,Z_j=O(j^{-1})$.

These estimates must also hold on histories not selected at a given
index.  Fix a $J$-history and $J/2\le j\le3J$.  Let $J'$ be the
dyadic parameter selected at $j$, so $J'\le j<2J'$.  Then $J$ and
$J'$ are comparable, and $j$ belongs to fixed positive-fraction
windows of both histories.  Apply
Lemma~\ref{lem:trajectory-overlap} with the selected state as the
unprimed state.  Since that state has $u,Z=O(j^{-1})$, the two
$u$ coordinates differ by $O(j^{-2})$, and so do the two $Z$
coordinates.  Lemma~\ref{lem:bulk-odd} then gives the same remainder
bounds on the $J$-history.  Thus
\eqref{eq:trajectory-balanced-window} holds throughout every
window asserted in the proposition, with uniform constants.

Finally undo the balance transformations and the subtraction of the
analytic Gaussian curve used to define $u,B_e,B_o$.  The coordinate
relations in \eqref{eq:bulk-coordinate-conversion} give
\begin{align*}
 t&=u+O(u^2+Z^2),\qquad |z|\le C|Z|,\\
 \|\mathcal R\|
   &\le C\bigl(u^2+Z^2+\|B_e\|+\|B_o\|\bigr).
\end{align*}
The quadratic order in the last line includes the remainder of the
Gaussian curve, whose constant and linear gradient terms have
already been separated.  Applying these relations to
\eqref{eq:trajectory-balanced-window} proves
\eqref{eq:trajectory-bulk-window}.
\end{proof}

The argument uses uniform convergence on positive-fraction windows
and the common physical readouts with their stated nonlinear errors.
The additional errors from entry-scale smoothing
and cutoff removal are exponentially small on those windows.
No rate of convergence in the height scaling limit was used to
obtain the coefficient $1/(2\log L)$.
